% paper/sections/05b_ccd_extensions.tex
%
% §4 CCD 法の続き（04_ccd.tex からの分割）
% 内容：非一様格子拡張・2次元混合偏微分・楕円型ソルバーとしての双対的役割
%
\subsection{非一様格子への拡張：座標変換法}
\label{sec:ccd_metric}

\noindent\textbf{本節の位置づけ：}
§~\ref{sec:CCD}（CCD 法の定式化）で確立した等間隔格子上の 6 次精度演算子を，
本章の非一様格子（§~\ref{sec:grid_gen}）に拡張する．

本稿では界面近傍の解像度を上げるために非一様格子を用いるが，CCD 係数（式~\eqref{eq:coef_CCD}）は等間隔格子 $h$ を前提としている．
非一様格子上で高精度を維持するための標準的な手法は，計算空間への座標変換を用いることである．

\noindent\textbf{座標変換による実装．}
物理空間 $x$（非一様）を，計算空間 $\xi$（一様，$\Delta\xi=1$）に写像する．
微分計算はすべて計算空間 $\xi$ 上で標準の等間隔 CCD スキームを用いて行い，その結果を連鎖律（chain rule）で物理空間の微分値に変換する．

\textbf{変換公式：}
物理空間での一次微分 $\partial f/\partial x$ および二次微分 $\partial^2 f/\partial x^2$ は，メトリック係数 $J = \partial \xi / \partial x$ を用いて次のように計算される：
\begin{align}
  \pd{f}{x} &= \pd{\xi}{x}\pd{f}{\xi} = J \fp{\xi} \\
  \pdd{f}{x} &= J \pd{}{x}\!\left(\fp{\xi}\right) + \fp{\xi}\pd{J}{x} = J^2 \fpp{\xi} + J \pd{J}{\xi} \fp{\xi}
\end{align}
ここで，$f^{(\cdot)}_\xi$ は計算空間上で CCD 法により求めた微分値である．
この手法により，複雑な非一様格子係数を導出することなく，一様格子の CCD ソルバーをそのまま再利用できる利点がある．

\medskip

\begin{tcolorbox}[warnbox, title={$\partial J/\partial\xi$ の精度：全体精度への影響}]
式~\eqref{eq:transform_2nd_correct}（第~\ref{sec:grid}章で導出）の右辺第2項には $\partial J/\partial\xi$ が含まれる．
$J_x$ を素朴な $\Ord{h^2}$ 中心差分で評価した場合，
$\partial J/\partial\xi$ の精度が $\Ord{h}$ 以下となり，
CCD の $\Ord{h^6}$ 空間精度全体が $\Ord{h}$ または $\Ord{h^2}$ に劣化する．

\textbf{解決策：} 座標配列 $x(\xi)$ を「関数値」として CCD に入力し，
$dx/d\xi$ および $d^2x/d\xi^2$ を $\Ord{h^6}$ 精度で求める：
\[
  J|_i = \frac{1}{(dx/d\xi)_i}, \qquad
  \pd{J}{\xi}\bigg|_i = -\frac{(d^2x/d\xi^2)_i}{(dx/d\xi)_i^2}
\]
これにより $J$ および $\partial J/\partial\xi$ がともに $\Ord{h^6}$ 精度で評価され，
物理空間の2階微分も $\Ord{h^6}$ 精度が保たれる（\ref{box:grid_jx_accuracy}の注意ボックスも参照）．
\end{tcolorbox}

\subsection{2次元への拡張：混合偏微分 \texorpdfstring{$D_{xy}^{(11)}$}{Dxy\^{}(11)} の計算}
\label{sec:ccd_mixed}

曲率 $\kappa$ の計算式（式~\eqref{eq:curvature_2d}）には混合偏微分 $\phi_{xy} = \partial^2\phi/\partial x\partial y$ が不可欠である．
CCD 法は1次元のブロック Thomas 法として定式化されているため，混合微分は以下の\textbf{2段階適用}で計算する．

\subsubsection{基本的な考え方：逐次的な1次元 CCD の適用}

\begin{tcolorbox}[resultbox, title={混合微分 $\partial^2 f/\partial x\partial y$ の計算}]
\textbf{Step 1：$x$ 方向 CCD}

各行（$j = \text{const}$）に対して $x$ 方向のブロック Thomas 法を適用し，
各格子点 $(i,j)$ における $x$ 方向一次微分 $g_{i,j} \equiv \partial f/\partial x \big|_{(i,j)}$ を計算する．
この操作を全行（$j = 0, 1, \ldots, N_y-1$）に対して実行する．
\[
  g_{i,j} = D_x^{(1)} f \big|_{(i,j)} \quad \Leftarrow \quad x \text{ 方向 CCD（行スウィープ）}
\]

\textbf{Step 2：$g = \partial f/\partial x$ への $y$ 方向 CCD}

Step 1 で得た $g_{i,j}$ を，今度は\textbf{$y$ 方向の関数値}とみなして，
各列（$i = \text{const}$）に対して $y$ 方向のブロック Thomas 法を適用する．
\[
  \phi_{xy,\,i,j} = D_y^{(1)} g \big|_{(i,j)} = D_y^{(1)}\bigl(D_x^{(1)} f\bigr)\big|_{(i,j)}
\]

\textbf{結果：} 混合微分 $\partial^2 f/\partial x \partial y$ が全格子点で得られる．
\end{tcolorbox}

\subsubsection{精度の議論：\texorpdfstring{$O(h^6)$}{O(h\textasciicircum 6)} の維持}

各方向の CCD が $\Ord{h^6}$ 精度を持つとき，Step 2 での入力誤差増幅により一般には $\Ord{h^5}$ に劣化しうる．
ただし，$f \in C^8$ のもとで打ち切り誤差場 $E_g \propto h^6 f^{(7)}$ の $y$ 方向変動が $\Ord{h^6}$ に収まるため $\Ord{h^6}$ が維持される
（詳細は付録~\ref{app:ccd_mixed_precision}）．
格子収束テスト（第~\ref{sec:verification}章）で $N$ を倍にするたびに誤差が $1/64$ 倍（6次収束）かを確認すること．

\medskip

% 実装上の注意（メモリ管理・計算順序・境界条件の設定）は §~\ref{sec:algorithm} を参照．

\medskip
以上で CCD 法に基づく空間離散化スキームの構築が完了した．
次章では，これらの空間演算子を時間軸上で整合的に前進させる
時間積分フレームワーク（TVD-RK3 および AB2+IPC+Crank--Nicolson）を設計する．

